\chapter{Problem Formulation}
\label{ch:problem}

\section{The prediction and control problem, and the tensor format used to describe it}
\label{sec:forecasting_tensor}

This chapter defines the problem this thesis solves and the mathematical objects used to describe it.
The physical picture is longitudinal train dynamics (LTD): a chain of vehicles connected by couplers, moving along a route, pushed and slowed by traction, braking, resistance, grade, and the forces passed through the couplers themselves.
None of this physics is new. What is new is how it is packaged. This thesis represents the coupled-vehicle system as a small set of tensors, structured arrays, so that the simulator's outputs, the model's training targets, and the surrogate model built in later chapters all share the same layout, no matter how many vehicles are in the train.

The rest of this chapter proceeds as follows. \Cref{sec:prediction_task} states, in plain terms, what is being predicted and why it matters. \Cref{sec:classical_ltd_compact} restates the underlying physics in its usual form. \Cref{sec:tensor_state_rep,sec:tensor_dynamics_form} introduce the node and edge tensors and rewrite the dynamics in that format. \Cref{sec:forecasting_objective} states the learning objective, including the requirement that the learned model be differentiable so that it can later be used to choose control actions, not just predict outcomes. The reference simulator that generates training data from this formulation is documented separately, in \Cref{ch:refsim}.

\subsection{Physical setting and prediction task}
\label{sec:prediction_task}

Picture a freight train moving along a route. Before it runs, or while planning how it should run, a few questions matter: how much fuel will this trip use, will it arrive on time, and will the forces between cars ever get dangerously high. A more useful question is how the answers to those change if the driving plan is adjusted slightly.

Answering that last question well is the real goal of this thesis. A model that can only report what happened in one specific run is a forecasting tool. A model that can also say how the outcome would shift if traction or braking commands were changed, and can do so cheaply and repeatedly, becomes something a control algorithm can use to choose better driving plans on its own. That is the target: a model fast and smooth enough to sit inside an optimization loop that picks fuel-efficient throttle and brake commands, while staying within speed limits and coupler-force safety limits.

To make this concrete, every scenario is described by three things: a route profile (the track's grade, curvature, and speed limits along its length), a consist configuration (the ordered list of vehicles, their masses, resistance properties, and coupler characteristics), and a control or operating plan (the traction and braking commands applied over time, possibly adjusted by a driver or automation model). Given these three inputs, the quantities of interest include:

\begin{itemize}
    \item \textbf{Motion:} the position, velocity, and acceleration of every vehicle.
    \item \textbf{In-train forces:} the force at each coupler, which determines whether neighboring cars are pulling apart (draft) or pushing together (buff).
    \item \textbf{Energy and fuel:} the work or power used over the run.
    \item \textbf{Schedule:} trip time and arrival at scheduled points.
    \item \textbf{Safety margin:} how close any of the above comes to an operating limit.
\end{itemize}

This thesis treats the physics simulator described in \Cref{ch:refsim} as ground truth when its parameters are fixed. Uncertainty in those parameters, or in the environment, is folded in as well (\cref{sec:uncertainty_sources}), but its role here is narrower than a full probabilistic forecast: it mainly shows up as a safety margin that the control algorithm has to respect, rather than as a result reported for its own sake.

\subsection{Classical longitudinal dynamics formulation}
\label{sec:classical_ltd_compact}

Classical longitudinal modeling represents the train as a chain of $N$ vehicles indexed $i=0,\ldots,N-1$ from lead to tail, connected by $N-1$ couplers.
Let $x_i(t)$ and $v_i(t)=\dot{x}_i$ denote the longitudinal position and velocity of vehicle $i$ along the route, and let $m_i$ denote its mass.
Between vehicles $i$ and $i+1$, coupler $i$ is associated with a kinematic extension $\delta_i$ (relative displacement from nominal spacing, positive in draft in the convention used below) and a coupler force $F_i$ transmitted to the trailing vehicle in the pair.
Resistance, grade, and optional curvature terms enter as external longitudinal forces on each vehicle; traction and braking appear as controlled inputs, possibly after actuator dynamics.
Summarizing the external and control contributions on vehicle $i$ by $F_{\mathrm{ext},i}(t)$ and the coupler incidence by the force on the rear of $i$ minus the force on the front, the translational dynamics take the familiar form
\begin{equation}
    \dot{x}_i = v_i, \qquad
    m_i \dot{v}_i = F_{\mathrm{ext},i}(t) + F_{i-1}(t) - F_i(t),
    \label{eq:classical_ltd_compact}
\end{equation}
with $F_{-1}=F_{N-1}=0$ when end forces are absent.
The coupler forces $F_i$ are determined by constitutive laws $F_i = \phi_i(\delta_i,\dot{\delta}_i,\ldots)$ linking kinematics to internal forces; slack, asymmetric buff and draft stiffness, and damping enter here (\cref{ch:refsim}).
When actuator lag or power limits are modeled, \cref{eq:classical_ltd_compact} is augmented by additional state variables per vehicle (for example brake and traction command states) driven by exogenous commands.
This is the standard description of a coupled multi-body train, and it holds regardless of the tensor notation introduced next.

\subsection{Tensorized state representation}
\label{sec:tensor_state_rep}

For learning, and for software that must handle variable train length and rich per-vehicle features, it is convenient to organize the same physics into two matrices that make the vehicle-coupler graph explicit.
In this thesis, a \emph{node} (vehicle) tensor $\mathbf{H}(t) \in \mathbb{R}^{N \times d_n}$ collects $d_n$ channels associated with each vehicle at time $t$.
Typical channels include longitudinal position and velocity, internal actuator or command states when those are part of the dynamic model, and static attributes replicated on each row (mass, Davis coefficients, traction capability flags, force limits, and similar parameters) so that a single array suffices for batched computation and for models that read local vehicle features.
An \emph{edge} (coupler) tensor $\mathbf{E}(t) \in \mathbb{R}^{(N-1) \times d_e}$ holds $d_e$ channels per coupler, such as extension $\delta_i$, relative speed $\dot{\delta}_i$, algebraic coupler force $F_i$, nominal spacing, slack half-width, and stiffness and damping parameters for buff and draft segments.
The assignment of channels is fixed in the reference implementation; the important point is the separation between vehicle-centered quantities in $\mathbf{H}$ and coupler-centered quantities in $\mathbf{E}$.

Along a discrete time grid $t_0,\ldots,t_{T-1}$, simulated trajectories can be stored as a \emph{rollout tensor} $\mathbf{X} \in \mathbb{R}^{T \times N \times d}$, whose time slice $\mathbf{X}_{k,:,:}$ is the node feature matrix at $t_k$ (so one may take $d=d_n$ and identify $\mathbf{X}_{k,:,:}$ with $\mathbf{H}(t_k)$ after any bookkeeping for channels).
Coupler histories are stored analogously as $\mathbf{E}_{\mathrm{hist}} \in \mathbb{R}^{T \times (N-1) \times d_e}$.

This layout keeps the structure of the physical problem visible: each coupler row depends only on the two vehicle rows next to it, mirroring the chain shape of \cref{eq:classical_ltd_compact}.
That chain shape is also why this thesis treats the train as a graph, with vehicles as nodes and couplers as edges, rather than as a fixed-length sequence.
A graph-based model can process trains of any length directly, passing information between neighboring vehicles one coupler at a time, without needing to pad shorter trains or mask out unused positions the way sequence models do.
Handling trains of different lengths this way, as a direct consequence of the representation rather than as an engineering workaround, is central to the train-agnostic goal of this thesis: the same trained model should work on a train it has never seen, of a length it has never seen, without modification.

\subsection{Dynamics in tensor-oriented form}
\label{sec:tensor_dynamics_form}

The tensor view is a reorganization of state and parameters, not a change to the underlying force balances.
Let $\mathbf{H}_{\mathrm{dyn}}(t)$ denote the dynamic channels within $\mathbf{H}(t)$ (for example positions, velocities, and actuator states in the extended reference model), and let $\mathbf{p}$ collect fixed parameters, including static columns of $\mathbf{H}$ and of $\mathbf{E}$ where those encode coupler constants.
The differential part of the evolution can be written compactly as
\begin{equation}
    \frac{\mathrm{d}}{\mathrm{d}t}\,\mathrm{vec}\!\left(\mathbf{H}_{\mathrm{dyn}}(t)\right)
    = \mathcal{F}_n\bigl(\mathbf{H}(t), \mathbf{E}(t), \mathcal{U}(t), \mathcal{R}; t\bigr),
    \label{eq:tensor_node_evolution}
\end{equation}
where $\mathcal{U}$ denotes the control input signals, $\mathcal{R}$ the route profile, and $\mathcal{F}_n$ the assembled right-hand side obtained from \cref{eq:classical_ltd_compact} and the actuator relations.
Coupler kinematics link $\mathbf{H}_{\mathrm{dyn}}$ to the extension and rate channels in $\mathbf{E}(t)$; constitutive laws determine the force channels.
In the reference integrator used for data generation in this thesis, several edge channels (including $\delta_i$, $\dot{\delta}_i$, and $F_i$) are \emph{reconstructed algebraically} at each evaluation from vehicle positions and velocities and from coupler parameters, rather than integrated as separate states.
Thus the edge tensor is in part an explicit bookkeeping of quantities that satisfy relations of the form
\begin{equation}
    \mathbf{E}(t) = \mathcal{G}_e\bigl(\mathbf{H}(t), \mathbf{p}\bigr),
    \label{eq:tensor_edge_map}
\end{equation}
for the algebraic channels, while \cref{eq:tensor_node_evolution} carries the true differential degrees of freedom.
If additional internal coupler or draft-gear states were introduced in future work, \cref{eq:tensor_edge_map} would be supplemented by differential equations for the corresponding rows of $\mathbf{E}(t)$, i.e.\ a nontrivial $\mathcal{F}_e$ on those components alongside $\mathcal{F}_n$.

The surrogate model built later in this thesis does not learn $\mathcal{F}_n$, and it does not learn a correction to it.
It works one \emph{step} at a time rather than one derivative at a time: given the state at $t_k$, it predicts the state at $t_{k+1}=t_k+h$ directly, for a fixed step $h$ matched to the spacing at which trajectories are stored.
Writing $\mathbf{s}_k$ for the dynamic state at time $t_k$, the model takes the form
\begin{equation}
    \mathbf{s}_{k+1}
    = \mathcal{S}_\theta\bigl(\mathbf{s}_k, \mathbf{E}(t_k), \mathcal{U}, \mathcal{R};\, h\bigr),
    \label{eq:tensor_step_operator}
\end{equation}
in which \cref{eq:tensor_node_evolution} appears only as the relation the training \emph{data} satisfy, not as a term the model evaluates at run time.

Within that step the four dynamic channels are not treated alike, and the division of labour is the substance of the design.
Position obeys $\dot{x}_i = v_i$ exactly, so it is recovered kinematically from the velocities at the two ends of the step rather than learned.
The brake and traction lag states obey $\dot{z} = (u - z)/\tau$ with the command $u$ known in advance, which is integrable in closed form across the step; giving them to a network would be asking it to learn a first-order lag it can simply be told.
Only velocity is genuinely learned: the network predicts the \emph{mean} acceleration of each vehicle over the step, together with a small correction on the per-coupler extension, which the kinematic reconstruction cannot recover from endpoint velocities alone.

Two consequences are worth stating here, because they shape everything that follows.
The first is that two of the four dynamic channels leave the network entirely and a third is mostly imposed, so what is learned narrows to two scalars per vehicle rather than a correction on every channel: the physics in this model is structural, not additive.
The second is that the step is the unit of computation.
The reference integrator takes many internal steps between stored samples, and the surrogate replaces all of them with a single evaluation, which is where its speed advantage comes from and why that advantage can be quoted from the integrator's own configuration rather than from a hardware comparison.
\Cref{sec:arch_anchor,sec:arch_channels} give the measurements that fix which channel is anchored on what; the point here is only that the formulation is a \emph{learned step} on the tensor state, with physics imposed channel by channel.

\subsection{Forecasting objective and simulator operator}
\label{sec:forecasting_objective}

The learning problem is to predict, under computation and uncertainty constraints, the outputs described in \cref{sec:prediction_task} from the structured inputs $(\mathcal{R}, \mathcal{C}, \mathcal{U})$.
The reference physics implements a deterministic operator
\begin{equation}
    F : (\mathcal{R}, \mathcal{C}, \mathcal{U}) \mapsto \mathbf{y},
    \label{eq:simulator_operator}
\end{equation}
where $\mathbf{y}$ collects trajectories (for example stacks of $\mathbf{H}(t_k)$ and $\mathbf{E}(t_k)$), scalar summaries such as energy and trip time, and any derived risk indicators.
The triple $(\mathcal{R}, \mathcal{C}, \mathcal{U})$ denotes the same structured inputs as the route profile $\route$, consist $\consist$, and control plan $\controlplan$ detailed in \cref{sec:input_definitions}.
When uncertain parameters and environment are modeled by a random vector $\xi \sim p(\xi)$, the same operator is written $F(\mathcal{R}, \mathcal{C}, \mathcal{U}; \xi)$ and outputs become random (\cref{sec:uncertainty_sources}).

The goal of this thesis is not just to match individual trajectories point by point.
The surrogate needs to be accurate enough that its outputs, especially coupler forces and trip energy, can be trusted for planning and for staying within safety margins.
Where uncertainty matters, it mainly shows up as a margin that the surrogate and the downstream control algorithm must respect, rather than as a full calibrated probability distribution reported over every output.

The surrogate introduced later in this thesis approximates $F$ by iterating the learned step of \cref{eq:tensor_step_operator}, using a message-passing graph neural network that operates directly on the node and edge tensors defined above (\cref{sec:tensor_dynamics_form}).
Because the surrogate must eventually be used to choose control actions, not just predict outcomes for a fixed plan, it is also built to be differentiable: it can be evaluated many times in a batch, and the gradient of its outputs with respect to the control plan can be computed directly.
The step formulation helps here rather than hindering.
A fixed update around a single network evaluation leaves no adaptive solver in the gradient path, so the sensitivity of a trajectory to the control plan is obtained by differentiating a known, fixed composition of steps.
That property is what later makes it possible to use the surrogate as the model inside a fuel-optimal control routine, instead of only as a fast forecasting tool.

\section{Input definitions}
\label{sec:input_definitions}

\subsection{Route profile}

A route profile $\route$ provides the exogenous fields sampled by each vehicle as it moves along the track.
In the reference simulator, $\route$ is a piecewise representation over distance $s$ that supports interpolation of grade, curvature, and speed constraints:
\begin{itemize}
    \item grade field $\theta(s)$ (or $\sin\theta(s)$ in implementation),
    \item curvature proxy $\kappa(s)$ when enabled,
    \item operational constraints such as speed limits $v_{\max}(s)$.
\end{itemize}
Given vehicle positions $x_i(t)$ from the node tensor, route-dependent forces are evaluated per vehicle at $s=x_i(t)$.
This keeps the forcing consistent with the tensorized state layout while preserving the classical physics of \cref{sec:classical_ltd_compact}.

Route profiles used in this thesis are built from real elevation and track alignment data rather than invented from scratch, so that the grades and curves seen during training and evaluation resemble those of an actual rail corridor.
Extra scenarios are produced by mildly perturbing these real profiles, for example in grade severity or adhesion, rather than by generating routes with no real-world basis.
The pipeline that turns raw map and elevation data into these profiles is documented in \cref{ch:routes}.

\subsection{Consist configuration}

A consist configuration $\consist$ defines the ordered vehicle chain and coupler parameters.
It is represented as two aligned tables that map directly to tensor channels:
\begin{itemize}
    \item \textbf{Vehicle attributes} for $i=0,\ldots,N-1$: mass, Davis coefficients, traction/brake limits, and actuator constants.
    \item \textbf{Coupler attributes} for $j=0,\ldots,N-2$: nominal spacing, slack half-width, buff/draft stiffness, and damping parameters.
\end{itemize}
Static vehicle attributes are stored as fixed channels of the node tensor $\mathbf{H}$, and static coupler attributes as fixed channels of the edge tensor $\mathbf{E}$.
This organization is exactly the data schema used by the simulator and by the surrogate model.

\subsection{Control plan}

A control plan $\controlplan$ specifies exogenous command signals over time (or equivalent distance-indexed schedules mapped to time during simulation).
For each powered or braked vehicle, commands are represented as channels $u_{\mathrm{trac},i}(t)$ and $u_{\mathrm{brk},i}(t)$ driving the actuator dynamics documented in \cref{ch:refsim}.
The simulator therefore receives structured inputs as the triplet $(\route,\consist,\controlplan)$ introduced in \cref{eq:simulator_operator}, and produces tensorized trajectories and derived summaries.

\section{Output definitions}

The output of one simulation run is a structured object
\begin{equation}
    \mathbf{y} = \Bigl(\mathbf{X},\,\mathbf{E}_{\mathrm{hist}},\,\mathbf{s},\,\mathbf{r}\Bigr),
\end{equation}
where $\mathbf{X}$ and $\mathbf{E}_{\mathrm{hist}}$ are rollout tensors, $\mathbf{s}$ contains scalar summaries, and $\mathbf{r}$ contains risk-oriented statistics.

\subsection{Trajectory-level outputs}

At discrete times $t_k$, the node rollout is $\mathbf{X}_{k,:,:}\in\mathbb{R}^{N\times d_n}$ and the edge rollout is $\mathbf{E}_{\mathrm{hist},k,:,:}\in\mathbb{R}^{(N-1)\times d_e}$.
These tensors contain the channels defined in \cref{sec:tensor_state_rep}: per-vehicle motion and actuator states in $\mathbf{X}$, and per-coupler extension, relative speed, and force channels in $\mathbf{E}_{\mathrm{hist}}$.
Classical quantities are recovered directly from selected channels, e.g.\ $x_i(t_k)$, $v_i(t_k)$, and coupler force $F_j(t_k)$.
For notation, the coupler force channel may be written as
\begin{equation}
    F_j(t) = f_j\!\bigl(\route,\consist,\controlplan,t\bigr),
    \label{eq:coupler_force}
\end{equation}
where $f_j$ is induced by the simulator dynamics and constitutive laws.

\subsection{Aggregate outputs}

Aggregate outputs are deterministic functionals of rollout tensors for fixed parameters.
Typical summaries used throughout this thesis include
\begin{align}
    E_{\mathrm{trip}} &= \int_0^{T} P_{\mathrm{trac}}(t)\,\mathrm{d}t, \label{eq:energy} \\
    T_{\mathrm{arr}} &= \inf\{t:\;x_{\mathrm{lead}}(t)\ge s_{\mathrm{target}}\}, \label{eq:arrival_time} \\
    F_{\max} &= \max_{t,j}\,|F_j(t)|.
\end{align}
All such metrics are post-processing maps on $(\mathbf{X},\mathbf{E}_{\mathrm{hist}})$ and therefore remain consistent with the tensorized representation.

\subsection{Risk metrics}
\label{sec:risk_metrics}

For uncertainty-aware forecasting, risk is expressed as probabilities and quantiles of tensor-derived observables.
For threshold $\tau$ and quantity $Z$ (for example $F_{\max}$ or $|F_j(t)|$), the exceedance probability is
\begin{equation}
    \mathbb{P}(Z>\tau).
    \label{eq:exceedance_prob}
\end{equation}
Likewise, predictive quantiles satisfy $\mathbb{P}(Y\le Q_p)=p$ for selected levels $p$.
Both quantities are evaluated on variables computed from node/edge rollouts, so no separate non-tensor output definition is required.
In this thesis these risk metrics are used mainly for coupler-force safety, as discussed next.

\section{Uncertainty sources}
\label{sec:uncertainty_sources}

Not every input to the simulator is known exactly.
Vehicle masses, resistance coefficients, coupler stiffness, the exact shape of the route, and even how precisely a command is executed all carry some uncertainty.
This section lists where that uncertainty comes from and how this thesis handles it.
Because the main use of the surrogate in this thesis is fuel-optimal control, not general-purpose forecasting, uncertainty is treated mainly as something that narrows the safe operating margin the control algorithm must respect, rather than as a full predictive distribution reported for its own sake.

Uncertainty is introduced through a random vector $\xi$ that perturbs physically meaningful components of the structured input and model:
\begin{enumerate}
    \item \textbf{Vehicle and coupler parameters:} mass, resistance coefficients, slack and stiffness values, brake/traction effectiveness.
    \item \textbf{Route and environment:} grade/curvature profile errors, adhesion-related effects, and external disturbances.
    \item \textbf{Operational realization:} command timing and execution variability.
    \item \textbf{Model discrepancy:} residual mismatch between the reference simulator and real operations.
\end{enumerate}

\subsection{Formal treatment of uncertainty}

Let $\xi\sim p(\xi)$ denote the joint uncertainty law.
Then the simulator operator from \cref{eq:simulator_operator} becomes
\begin{equation}
    \mathbf{y}(\xi)=F(\route,\consist,\controlplan;\xi),
    \label{eq:probabilistic_simulator}
\end{equation}
and every tensor channel and derived metric becomes a random variable.
For example, the exceedance probability for coupler interface $j$ is
\begin{equation}
    \mathbb{P}_{\xi}\!\bigl(F_j(t;\xi)>\tau\bigr)
    =\int_{\{\xi:\,F_j(t;\xi)>\tau\}}p(\xi)\,\mathrm{d}\xi.
    \label{eq:exceedance_prob_formal}
\end{equation}
In practice, these probabilities are approximated by Monte Carlo simulation or by probabilistic surrogates trained on ensembles generated from the same tensorized simulator.
The role this plays in the thesis is narrow: $\mathbb{P}_\xi(F_j > \tau)$ is used to define a safety margin that the fuel-optimal control routine is required to respect, for example by keeping the estimated coupler force below its limit with high probability, rather than as a general-purpose calibrated forecast reported across every output variable.

\section{Supervised Learning Objective}

Training data are generated by the reference simulator documented in \cref{ch:refsim}:
\begin{equation}
    \mathcal{D}=\left\{
    \bigl(\route^{(k)},\consist^{(k)},\controlplan^{(k)},\mathbf{X}^{(k)},\mathbf{E}_{\mathrm{hist}}^{(k)},\mathbf{s}^{(k)}\bigr)
    \right\}_{k=1}^{M}.
\end{equation}
The surrogate $f_\theta$ maps structured inputs to trajectory and summary predictions in the same tensor layout, by iterating the learned step of \cref{eq:tensor_step_operator}.
A generic training objective is
\begin{equation}
    \mathcal{L}(\theta)=
    \lambda_X \mathcal{L}_X +
    \lambda_E \mathcal{L}_E +
    \lambda_S \mathcal{L}_S,
    \label{eq:training_loss}
\end{equation}
with losses on the node channels, on the edge channels, and on the trip-level summaries.
Supervision is on the learned quantities only: the imposed channels of \cref{sec:tensor_dynamics_form} are reconstructed exactly and contribute no error to fit.

\subsection{Loss term definitions}

\textbf{Rollout losses.} $\mathcal{L}_X$ and $\mathcal{L}_E$ compare predicted and reference tensor channels over time, vehicle/coupler index, and feature channels (with masking for variable train length where needed).

\textbf{Summary loss.} $\mathcal{L}_S$ penalizes errors on scalar outputs such as trip energy, arrival time, or peak coupler force.

\textbf{On physical consistency.} No separate physics-consistency penalty appears in \cref{eq:training_loss}, and none is needed.
Under the channel-wise construction of \cref{sec:tensor_dynamics_form} the kinematic relation between position and velocity, and the actuator lag dynamics, hold by construction at every step rather than approximately and on average.
A penalty term is the right tool when a model is free to violate a relation and must be discouraged from doing so; here the model cannot express the violation in the first place.

\textbf{On the safety margin.} The coupler-force margin that the control routine relies on (\cref{sec:risk_metrics}) is treated as an evaluation target rather than as a separate distributional head.
This follows the narrower role assigned to uncertainty in \cref{sec:uncertainty_sources}: what the controller needs is a trustworthy estimate of the peak force a plan will produce, and that is a question about the accuracy of the predicted force channel, not about calibrating a distribution over it.

\textbf{Remark.} \Cref{sec:arch_objective} instantiates these terms for the graph neural network used in this thesis, and the training chapter adds the perturbation scheme that makes a single-step model stable when iterated over long horizons.

\section{Conceptual pipeline}

The full workflow, from generating data to using the trained model for control, follows the same tensor objects at every step:

\begin{enumerate}
    \item \textbf{Structured inputs}: Specify $(\route,\consist,\controlplan)$ and, for uncertainty studies, sample $\xi$.
    \item \textbf{Reference simulation}: Run the tensorized LTD simulator (\cref{ch:refsim}) to obtain $(\mathbf{X},\mathbf{E}_{\mathrm{hist}},\mathbf{s})$.
    \item \textbf{Dataset assembly}: Store many such runs as supervised pairs for training and validation.
    \item \textbf{Surrogate training}: Fit $f_\theta$ to predict rollout tensors and derived summaries, including the coupler-force safety margin.
    \item \textbf{Forecasting}: Use $f_\theta$ for fast scenario evaluation in the same output space as the simulator.
    \item \textbf{Control}: Use $f_\theta$ and its gradients inside an optimization routine, such as model-predictive control, to choose traction and braking commands that reduce fuel use while respecting speed limits and coupler-force safety margins.
\end{enumerate}

\begin{figure}[h]
\centering
\caption{Conceptual pipeline from structured inputs to tensorized simulator outputs, surrogate forecasting, and fuel-optimal control.}
\label{fig:problem_pipeline}
\end{figure}

\section{Evaluation targets}

Evaluation is defined on the same outputs as the training objective:

\textbf{Trajectory fidelity.} How closely predicted vehicle motion and coupler forces track the reference simulator, channel by channel (for example MAE/RMSE on velocity and coupler-force channels).

\textbf{Summary accuracy.} Errors on trip-level numbers such as total energy used, arrival time, and peak coupler force.

\textbf{Safety-margin quality.} How well the surrogate's coupler-force margin estimates hold up, since this is what the control routine relies on to stay within limits.

\textbf{Generalization across trains and routes.} Accuracy on consist lengths and route corridors not seen during training. Since the surrogate is meant to work on trains it was never explicitly trained on, this is not a side check; it is one of the central claims of the thesis.

\textbf{Operational utility.} Inference-time speedup relative to the reference simulator, and performance as the model inside the fuel-optimal control routine, compared against a simple constant-speed driving plan.

\section{Chapter summary}

This chapter defined the prediction and control problem this thesis addresses.
Classical longitudinal dynamics provides the physical foundation (\cref{eq:classical_ltd_compact}).
This thesis recasts that same system into node and edge tensors $\mathbf{H}(t)$ and $\mathbf{E}(t)$ (\cref{sec:tensor_state_rep,sec:tensor_dynamics_form}), a graph-shaped representation that lets a single model handle trains of any length.
The simulator $F$ maps structured inputs $(\route, \consist, \controlplan)$ to structured outputs $\mathbf{y}$: trajectory tensors, trip-level summaries such as energy and arrival time, and safety-relevant quantities such as peak coupler force (\cref{eq:simulator_operator}).
We distinguished deterministic outputs under fixed parameters from outputs under random parameter uncertainty $\xi$ (\cref{eq:probabilistic_simulator}).

The learning task is hard for four reasons: the inputs are high-dimensional and structured, the outputs span different scales and types, physical consistency has to be preserved, and safety-relevant quantities need a notion of margin rather than a single number.
This thesis addresses the last of these narrowly, folding uncertainty into a coupler-force safety margin rather than building a full calibrated forecast of every output (\cref{sec:uncertainty_sources}).

The next chapter documents the reference simulator that generates the data used throughout the rest of this thesis.
\Cref{ch:architecture} introduces the surrogate model itself: a graph neural network that advances the tensor state one fixed step at a time, with position and actuator dynamics imposed rather than learned, and that is differentiable end to end, so that it can later serve as the model inside a fuel-optimal control routine.
